% Séquence 4 : Géométrie du triangle (partie 1) - 5e
\setseqtitle{Géométrie du triangle (partie 1)}
\chapter{Géométrie du triangle (partie 1)}

\begin{objectifsbox}
	\textbf{Objectifs d'apprentissage.} À l'issue de la séquence, l'élève sera capable de :
	\begin{itemize}
		\item Construire un triangle à partir de données suffisantes
		\item Déterminer si trois longueurs peuvent former un triangle
		\item Appliquer l'inégalité triangulaire
		\item Reconnaître et construire les hauteurs d'un triangle
		\item Reconnaître et construire les médianes d'un triangle
		\item Reconnaître et construire les médiatrices des côtés d'un triangle
		\item Connaître les points remarquables du triangle (orthocentre, centre de gravité, centre du cercle circonscrit)
	\end{itemize}
\end{objectifsbox}

\section{Rappels : Construction de triangles}

\subsection{Les instruments de géométrie}

\begin{remarkbox}
	Pour construire un triangle, on utilise principalement :
	\begin{itemize}
		\item \textbf{La règle graduée} : pour mesurer et tracer des segments
		\item \textbf{Le compas} : pour tracer des cercles et reporter des longueurs
		\item \textbf{Le rapporteur} : pour mesurer et tracer des angles
		\item \textbf{L'équerre} : pour tracer des angles droits et des perpendiculaires
	\end{itemize}
\end{remarkbox}

\subsection{Construction avec trois longueurs}

\begin{methodebox}[Construire un triangle connaissant trois longueurs]
	Pour construire le triangle $ABC$ tel que $AB = c$, $AC = b$ et $BC = a$ :
	\begin{enumerate}
		\item Tracer le segment $[AB]$ de longueur \solution[1.5cm]{$c$}.
		\item Tracer un arc de cercle de centre \solution[1cm]{$A$} et de rayon \solution[1.5cm]{$b$}.
		\item Tracer un arc de cercle de centre \solution[1cm]{$B$} et de rayon \solution[1.5cm]{$a$}.
		\item Les deux arcs se coupent en un point \solution[1cm]{$C$}.
		\item Tracer les segments $[AC]$ et $[BC]$.
	\end{enumerate}

	\medskip

	\textbf{Illustration des étapes :}

	\begin{center}
		\begin{tikzpicture}[scale=0.45]
			% Étape 1 : Segment AB
			\begin{scope}[xshift=0cm]
				\coordinate (A) at (0,0);
				\coordinate (B) at (5,0);
				\draw[thick, mainBlue] (A) -- (B);
				\fill (A) circle (2.5pt) node[below] {$A$};
				\fill (B) circle (2.5pt) node[below] {$B$};
				\node at (2.5,-1.2) {\small \textbf{Étape 1}};
				\node[mainBlue] at (2.5,0.4) {\small $c$};
			\end{scope}

			% Étape 2 : Arc de centre A
			\begin{scope}[xshift=7cm]
				\coordinate (A) at (0,0);
				\coordinate (B) at (5,0);
				\draw[thick] (A) -- (B);
				\draw[thick, mainRed] (A) ++(20:3.5) arc (20:140:3.5);
				\draw[mainRed, ->] (A) -- ++(80:3.5) node[midway, above left] {\small $b$};
				\fill (A) circle (2.5pt) node[below] {$A$};
				\fill (B) circle (2.5pt) node[below] {$B$};
				\node at (2.5,-1.2) {\small \textbf{Étape 2}};
			\end{scope}

			% Étape 3 : Arc de centre B
			\begin{scope}[xshift=14cm]
				\coordinate (A) at (0,0);
				\coordinate (B) at (5,0);
				\draw[thick] (A) -- (B);
				\draw[thick, mainRed] (A) ++(20:3.5) arc (20:140:3.5);
				\draw[thick, mainGreen] (B) ++(60:3) arc (60:160:3);
				\draw[mainGreen, ->] (B) -- ++(110:3) node[midway, above right] {\small $a$};
				\fill (A) circle (2.5pt) node[below] {$A$};
				\fill (B) circle (2.5pt) node[below] {$B$};
				\node at (2.5,-1.2) {\small \textbf{Étape 3}};
			\end{scope}

			% Étape 4 : Point C
			\begin{scope}[xshift=21cm]
				\coordinate (A) at (0,0);
				\coordinate (B) at (5,0);
				\coordinate (C) at (2.825,2.07);
				\draw[thick] (A) -- (B);
				\draw[dashed, mainRed!50] (A) ++(20:3.5) arc (20:140:3.5);
				\draw[dashed, mainGreen!50] (B) ++(60:3) arc (60:160:3);
				\fill (A) circle (2.5pt) node[below] {$A$};
				\fill (B) circle (2.5pt) node[below] {$B$};
				\fill[mainBlue] (C) circle (3pt) node[above] {$C$};
				\node at (2.5,-1.2) {\small \textbf{Étape 4}};
			\end{scope}

			% Étape 5 : Triangle complet
			\begin{scope}[xshift=28cm]
				\coordinate (A) at (0,0);
				\coordinate (B) at (5,0);
				\coordinate (C) at (2.825,2.07);
				\draw[thick, mainBlue] (A) -- (B) -- (C) -- cycle;
				\fill (A) circle (2.5pt) node[below] {$A$};
				\fill (B) circle (2.5pt) node[below] {$B$};
				\fill (C) circle (2.5pt) node[above] {$C$};
				\node at (2.5,-1.2) {\small \textbf{Étape 5}};
			\end{scope}
		\end{tikzpicture}
	\end{center}
\end{methodebox}

\begin{examplebox}
	\textbf{Construire le triangle $DEF$ tel que $DE = 6$ cm, $EF = 5$ cm et $DF = 4$ cm.}

	\medskip

	\ifcorrige
	\begin{center}
		\begin{tikzpicture}[scale=0.8]
			% Construction
			\coordinate (D) at (0,0);
			\coordinate (E) at (6,0);
			\coordinate (F) at (2.5,3.3);

			% Arcs de construction (en pointillés)
			\draw[dashed, mainBlue!50] (D) circle (4);
			\draw[dashed, mainRed!50] (E) circle (5);

			% Triangle
			\draw[thick] (D) -- (E) -- (F) -- cycle;

			% Points
			\fill (D) circle (2pt) node[below left] {$D$};
			\fill (E) circle (2pt) node[below right] {$E$};
			\fill (F) circle (2pt) node[above] {$F$};

			% Mesures
			\draw[<->, mainGreen] (D) -- node[below, fill=white] {6 cm} (E);
		\end{tikzpicture}
	\end{center}
	\fi
\end{examplebox}

\subsection{Construction avec deux longueurs et un angle}

\begin{methodebox}[Construire un triangle connaissant deux longueurs et l'angle compris]
	Pour construire le triangle $ABC$ tel que $AB = c$, $AC = b$ et $\widehat{BAC} = \alpha$ :
	\begin{enumerate}
		\item Tracer le segment $[AB]$ de longueur \solution[1.5cm]{$c$}.
		\item Placer l'angle $\widehat{BAC} = \alpha$ avec le \solution[3cm]{rapporteur}.
		\item Sur la demi-droite $[Ax)$, placer le point $C$ tel que $AC = \solution[1.5cm]{$b$}$.
		\item Tracer le segment $[BC]$.
	\end{enumerate}

	\medskip

	\textbf{Illustration des étapes :}

	\begin{center}
		\begin{tikzpicture}[scale=0.5]
			% Étape 1 : Segment AB
			\begin{scope}[xshift=0cm]
				\coordinate (A) at (0,0);
				\coordinate (B) at (5,0);
				\draw[thick, mainBlue] (A) -- (B);
				\fill (A) circle (2.5pt) node[below] {$A$};
				\fill (B) circle (2.5pt) node[below] {$B$};
				\node at (2.5,-1) {\small \textbf{Étape 1}};
				\node[mainBlue] at (2.5,0.4) {\small $c$};
			\end{scope}

			% Étape 2 : Angle alpha
			\begin{scope}[xshift=8cm]
				\coordinate (A) at (0,0);
				\coordinate (B) at (5,0);
				\coordinate (X) at (2.5,4.33);
				\draw[thick] (A) -- (B);
				\draw[thick, mainRed, ->] (A) -- (X);
				\draw[mainRed] (1,0) arc (0:60:1);
				\node[mainRed] at (1.3,0.5) {\small $\alpha$};
				\fill (A) circle (2.5pt) node[below] {$A$};
				\fill (B) circle (2.5pt) node[below] {$B$};
				\node at (2.5,-1) {\small \textbf{Étape 2}};
			\end{scope}

			% Étape 3 : Point C sur la demi-droite
			\begin{scope}[xshift=16cm]
				\coordinate (A) at (0,0);
				\coordinate (B) at (5,0);
				\coordinate (C) at (2,3.46);
				\draw[thick] (A) -- (B);
				\draw[thick, mainRed] (A) -- (C);
				\draw[dashed, mainRed!50] (C) -- (2.5,4.33);
				\draw[mainRed] (1,0) arc (0:60:1);
				\fill (A) circle (2.5pt) node[below] {$A$};
				\fill (B) circle (2.5pt) node[below] {$B$};
				\fill[mainBlue] (C) circle (3pt) node[above left] {$C$};
				\node[mainRed] at (1,2) {\small $b$};
				\node at (2.5,-1) {\small \textbf{Étape 3}};
			\end{scope}

			% Étape 4 : Triangle complet
			\begin{scope}[xshift=24cm]
				\coordinate (A) at (0,0);
				\coordinate (B) at (5,0);
				\coordinate (C) at (2,3.46);
				\draw[thick, mainBlue] (A) -- (B) -- (C) -- cycle;
				\fill (A) circle (2.5pt) node[below] {$A$};
				\fill (B) circle (2.5pt) node[below] {$B$};
				\fill (C) circle (2.5pt) node[above] {$C$};
				\node at (2.5,-1) {\small \textbf{Étape 4}};
			\end{scope}
		\end{tikzpicture}
	\end{center}
\end{methodebox}

\begin{examplebox}
	\textbf{Construire le triangle $ABC$ tel que $AB = 5$ cm, $AC = 4$ cm et $\widehat{BAC} = 60^\circ$.}

	\medskip

	\ifcorrige
	\begin{center}
		\begin{tikzpicture}[scale=0.9]
			% Triangle
			\coordinate (A) at (0,0);
			\coordinate (B) at (5,0);
			\coordinate (C) at (2,3.46);

			% Arc d'angle
			\draw[mainRed] (0.8,0) arc (0:60:0.8);
			\node[mainRed] at (1.2,0.4) {$60^\circ$};

			% Triangle
			\draw[thick] (A) -- (B) -- (C) -- cycle;

			% Points
			\fill (A) circle (2pt) node[below left] {$A$};
			\fill (B) circle (2pt) node[below right] {$B$};
			\fill (C) circle (2pt) node[above] {$C$};

			% Mesures
			\draw[<->, mainGreen] (A) -- node[below, fill=white] {5 cm} (B);
			\draw[<->, mainBlue] (A) -- node[left, fill=white] {4 cm} (C);
		\end{tikzpicture}
	\end{center}
	\fi
\end{examplebox}

\subsection{Construction avec une longueur et deux angles}

\begin{methodebox}[Construire un triangle connaissant un côté et deux angles]
	Pour construire le triangle $ABC$ tel que $AB = c$, $\widehat{CAB} = \alpha$ et $\widehat{ABC} = \beta$ :
	\begin{enumerate}
		\item Tracer le segment $[AB]$ de longueur $\solution[1.5cm]{c}$.
		\item En $A$, construire l'angle $\widehat{CAB} = \solution[1.5cm]{\alpha}$ avec le rapporteur.
		\item En $B$, construire l'angle $\widehat{ABC} = \solution[1.5cm]{\beta}$ avec le rapporteur.
		\item Les deux demi-droites se coupent en un point $\solution[1cm]{C}$.
	\end{enumerate}

	\medskip

	\textbf{Illustration des étapes :}

	\begin{center}
		\begin{tikzpicture}[scale=0.5]
			% Étape 1 : Segment AB
			\begin{scope}[xshift=0cm]
				\coordinate (A) at (0,0);
				\coordinate (B) at (5,0);
				\draw[thick, mainBlue] (A) -- (B);
				\fill (A) circle (2.5pt) node[below] {$A$};
				\fill (B) circle (2.5pt) node[below] {$B$};
				\node at (2.5,-1) {\small \textbf{Étape 1}};
				\node[mainBlue] at (2.5,0.4) {\small $c$};
			\end{scope}

			% Étape 2 : Angle en A
			\begin{scope}[xshift=8cm]
				\coordinate (A) at (0,0);
				\coordinate (B) at (5,0);
				\coordinate (X) at ($(A)+(50:5)$);
				\draw[thick] (A) -- (B);
				\draw[thick, mainRed, ->] (A) -- (X);
				\draw[mainRed] (1,0) arc (0:50:1);
				\node[mainRed] at (1.3,0.4) {\small $\alpha$};
				\fill (A) circle (2.5pt) node[below] {$A$};
				\fill (B) circle (2.5pt) node[below] {$B$};
				\node at (2.5,-1) {\small \textbf{Étape 2}};
			\end{scope}

			% Étape 3 : Angle en B
			\begin{scope}[xshift=16cm]
				\coordinate (A) at (0,0);
				\coordinate (B) at (5,0);
				\coordinate (X1) at ($(A)+(50:5)$);
				\coordinate (X2) at ($(B)+(110:5)$);
				\draw[thick] (A) -- (B);
				\draw[thick, mainRed, ->] (A) -- (X1);
				\draw[thick, mainGreen, ->] (B) -- (X2);
				\draw[mainRed] (1,0) arc (0:50:1);
				\draw[mainGreen] (4,0) arc (180:110:1);
				\node[mainRed] at (1.3,0.4) {\small $\alpha$};
				\node[mainGreen] at (3.7,0.5) {\small $\beta$};
				\fill (A) circle (2.5pt) node[below] {$A$};
				\fill (B) circle (2.5pt) node[below] {$B$};
				\node at (2.5,-1) {\small \textbf{Étape 3}};
			\end{scope}

			% Étape 4 : Point C et triangle
			\begin{scope}[xshift=24cm]
				\coordinate (A) at (0,0);
				\coordinate (B) at (5,0);
				\coordinate (C) at (3.49,4.15);
				\draw[thick, mainBlue] (A) -- (B) -- (C) -- cycle;
				\fill (A) circle (2.5pt) node[below] {$A$};
				\fill (B) circle (2.5pt) node[below] {$B$};
				\fill (C) circle (2.5pt) node[above] {$C$};
				\node at (2.5,-1) {\small \textbf{Étape 4}};
			\end{scope}
		\end{tikzpicture}
	\end{center}
\end{methodebox}

\begin{remarkbox}
	\textbf{Rappel :} La somme des angles d'un triangle est toujours égale à \solution[2cm]{$180^\circ$}.

	Donc, si on connaît deux angles, on peut calculer le troisième : $\widehat{ACB} = 180^\circ - \alpha - \beta$
\end{remarkbox}

\section{Inégalité triangulaire}

\subsection{Découverte}

\begin{activitybox}[Activité : Peut-on toujours construire un triangle ?]
	\textbf{Matériel :} règle, compas

	\medskip

	\textbf{Partie 1 : Essais de construction}

	Essayer de construire les triangles suivants en utilisant la méthode des trois longueurs :

	\begin{enumerate}[label=\alph*)]
		\item Triangle $ABC$ avec $AB = 6$ cm, $AC = 4$ cm et $BC = 5$ cm
		\item Triangle $DEF$ avec $DE = 7$ cm, $EF = 3$ cm et $DF = 2$ cm
		\item Triangle $GHI$ avec $GH = 5$ cm, $HI = 3$ cm et $GI = 8$ cm
	\end{enumerate}

	\medskip

	\textbf{Question 1 :} Quelles constructions ont réussi ? Lesquelles ont échoué ?

	\solution[8cm]{Les constructions a) réussit. Les constructions b) et c) échouent car les arcs de cercle ne se coupent pas.}

	\medskip

	\textbf{Partie 2 : Analyse}

	Compléter le tableau suivant pour chaque triangle :

	\begin{center}
		\begin{tabular}{|c|c|c|c|c|}
			\hline
			\textbf{Triangle} & \textbf{Côté 1} & \textbf{Côté 2} & \textbf{Côté 3} & \textbf{Construction possible ?} \\
			\hline
			$ABC$ & 6 cm & 4 cm & 5 cm & \solution[2.5cm]{Oui} \\
			\hline
			$DEF$ & 7 cm & 3 cm & 2 cm & \solution[2.5cm]{Non} \\
			\hline
			$GHI$ & 5 cm & 3 cm & 8 cm & \solution[2.5cm]{Non} \\
			\hline
		\end{tabular}
	\end{center}

	\medskip

	\textbf{Question 2 :} Pour chaque triangle, calculer la somme des deux plus petits côtés et comparer avec le troisième côté.

	\solution[10cm]{
		\begin{itemize}
			\item Triangle $ABC$ : $4 + 5 = 9 > 6$ donc construction possible
			\item Triangle $DEF$ : $3 + 2 = 5 < 7$ donc construction impossible
			\item Triangle $GHI$ : $3 + 5 = 8 = 8$ donc construction impossible (cas limite)
		\end{itemize}
	}

	\medskip

	\textbf{Question 3 :} Quelle condition doit vérifier les longueurs des côtés pour qu'un triangle puisse être construit ?

	\solution[10cm]{La somme des longueurs de deux côtés doit être strictement supérieure à la longueur du troisième côté.}
\end{activitybox}

\subsection{L'inégalité triangulaire}

\begin{proprietebox}[Inégalité triangulaire]
	Dans tout triangle $ABC$, la longueur d'un côté est toujours \textbf{\solution[4cm]{strictement inférieure}} à la somme des longueurs des deux autres côtés.

	On peut écrire :
	\begin{align*}
		AB &< \solution[1.5cm]{AC} + \solution[1.5cm]{BC}\\
		AC &< \solution[1.5cm]{AB} + \solution[1.5cm]{BC}\\
		BC &< \solution[1.5cm]{AB} + \solution[1.5cm]{AC}
	\end{align*}
\end{proprietebox}

\begin{remarkbox}
	\textbf{Interprétation :}

	L'inégalité triangulaire exprime que le chemin direct entre deux points est toujours plus court que le chemin qui passe par un troisième point.

	« Le plus court chemin entre deux points est la ligne droite. »
\end{remarkbox}

\begin{proprietebox}[Condition d'existence d'un triangle]
	Trois longueurs $a$, $b$ et $c$ peuvent former un triangle si et seulement si :
	$$a < b + c \quad \text{et} \quad b < a + c \quad \text{et} \quad c < a + b$$

	\textbf{En pratique :} Il suffit de vérifier que la \solution[4cm]{plus grande} longueur est strictement inférieure à la somme des deux autres.
\end{proprietebox}

\begin{examplebox}
	\textbf{Exemple 1 :} Les longueurs 5 cm, 7 cm et 10 cm peuvent-elles former un triangle ?

	\textbf{Solution :}

	La plus grande longueur est \solution[1.5cm]{10} cm.

	On vérifie : $5 + 7 = \solution[1.5cm]{12} \;\solution[0.8cm]{>} 10$

	Donc ces trois longueurs \solution[2cm]{peuvent} former un triangle.

	\medskip

	\textbf{Exemple 2 :} Les longueurs 3 cm, 4 cm et 8 cm peuvent-elles former un triangle ?

	\textbf{Solution :}

	La plus grande longueur est \solution[1.5cm]{8} cm.

	On vérifie : $3 + 4 = \solution[1.5cm]{7} \;\solution[0.8cm]{<} 8$

	Donc ces trois longueurs \solution[3.5cm]{ne peuvent pas} former un triangle.

	\medskip

	\textbf{Exemple 3 :} Les longueurs 4 cm, 5 cm et 9 cm peuvent-elles former un triangle ?

	\textbf{Solution :}

	La plus grande longueur est \solution[1.5cm]{9} cm.

	On vérifie : $4 + 5 = \solution[1.5cm]{9} \;\solution[0.8cm]{=} 9$

	Comme $4 + 5 = 9$ (et non $4 + 5 > 9$), ces trois longueurs \solution[3.5cm]{ne peuvent pas} former un triangle. Les trois points seraient alignés.
\end{examplebox}

\begin{methodebox}[Vérifier l'inégalité triangulaire]
	Pour vérifier si trois longueurs $a$, $b$, $c$ peuvent former un triangle :
	\begin{enumerate}
		\item Identifier la \solution[4cm]{plus grande} longueur.
		\item Calculer la \solution[2.5cm]{somme} des deux autres longueurs.
		\item Comparer : si la somme est \solution[5cm]{strictement supérieure} à la plus grande longueur, alors le triangle existe.
	\end{enumerate}
\end{methodebox}

\begin{applicationbox}[Problème]
	Un triangle $ABC$ a pour côtés $AB = 7$ cm et $BC = 5$ cm.

	\begin{enumerate}[label=\alph*)]
		\item Quelles sont les valeurs possibles (en cm) pour la longueur $AC$ ?

		\solution[10cm]{D'après l'inégalité triangulaire : $7 - 5 < AC < 7 + 5$, donc $2 < AC < 12$ cm.}

		\item Le troisième côté peut-il mesurer 2 cm ? 3 cm ? 12 cm ? 11 cm ?

		\solution[10cm]{
			\begin{itemize}
				\item 2 cm : Non (on doit avoir $AC > 2$)
				\item 3 cm : Oui ($2 < 3 < 12$)
				\item 12 cm : Non (on doit avoir $AC < 12$)
				\item 11 cm : Oui ($2 < 11 < 12$)
			\end{itemize}
		}
	\end{enumerate}
\end{applicationbox}

\section{Les droites remarquables du triangle}

\subsection{Les hauteurs}

\begin{definitionbox}[Hauteur d'un triangle]
	Dans un triangle, une \textbf{hauteur} est une droite qui passe par un \solution[3cm]{sommet} et qui est \solution[4cm]{perpendiculaire} au côté opposé (ou à son prolongement).

	Un triangle possède \solution[1.5cm]{trois} hauteurs.
\end{definitionbox}

\begin{examplebox}
	\textbf{Les trois hauteurs du triangle $ABC$}

	\ifcorrige
	\begin{center}
		\begin{tikzpicture}[scale=0.9]
			% Triangle
			\coordinate (A) at (0,0);
			\coordinate (B) at (6,0);
			\coordinate (C) at (2,4);

			% Pieds des hauteurs
			\coordinate (HA) at ($(B)!(A)!(C)$);
			\coordinate (HB) at ($(A)!(B)!(C)$);
			\coordinate (HC) at (2,0);

			% Triangle
			\draw[thick] (A) -- (B) -- (C) -- cycle;

			% Hauteurs
			\draw[mainRed, thick] (A) -- (HA);
			\draw[mainRed, thick] (B) -- (HB);
			\draw[mainRed, thick] (C) -- (HC);

			% Marques d'angles droits
			\draw (HC) ++(0,0.3) -- ++(0.3,0) -- ++(0,-0.3);
			\draw (HA) ++(-0.15,0.15) -- ++(0.21,0.21) -- ++(-0.21,-0.21);
			\draw (HB) ++(0.15,0.15) -- ++(-0.21,0.21) -- ++(0.21,-0.21);

			% Points
			\fill (A) circle (2pt) node[below left] {$A$};
			\fill (B) circle (2pt) node[below right] {$B$};
			\fill (C) circle (2pt) node[above] {$C$};

			% Légendes
			\node[mainRed] at (-0.8,2) {hauteur issue de $A$};
			\node[mainRed] at (6.8,2) {hauteur issue de $B$};
			\node[mainRed] at (2.5,1.5) {hauteur issue de $C$};
		\end{tikzpicture}
	\end{center}
	\fi

	\begin{itemize}
		\item La hauteur issue de $A$ est perpendiculaire à \solution[2cm]{$(BC)$}
		\item La hauteur issue de $B$ est perpendiculaire à \solution[2cm]{$(AC)$}
		\item La hauteur issue de $C$ est perpendiculaire à \solution[2cm]{$(AB)$}
	\end{itemize}
\end{examplebox}

\begin{remarkbox}
	\textbf{Cas particuliers :}
	\begin{itemize}
		\item Dans un triangle rectangle, deux hauteurs sont les côtés de l'angle droit.
		\item Dans un triangle obtusangle, deux hauteurs tombent à l'extérieur du triangle (sur le prolongement des côtés).
	\end{itemize}
\end{remarkbox}

\begin{proprietebox}[Orthocentre]
	Les trois hauteurs d'un triangle sont \solution[4cm]{concourantes}.

	Le point de concours des trois hauteurs s'appelle l'\textbf{\solution[4cm]{orthocentre}} du triangle.
\end{proprietebox}

\begin{methodebox}[Construire une hauteur]
	Pour construire la hauteur issue de $A$ dans le triangle $ABC$ :
	\begin{enumerate}
		\item Tracer la droite $(BC)$ (ou son prolongement si nécessaire).
		\item Avec l'\solution[2.5cm]{équerre} ou le compas, tracer la perpendiculaire à $(BC)$ passant par \solution[1.5cm]{$A$}.
		\item Cette droite est la hauteur issue de \solution[1.5cm]{$A$}.
	\end{enumerate}
\end{methodebox}

\subsection{Les médianes}

\begin{definitionbox}[Médiane d'un triangle]
	Dans un triangle, une \textbf{médiane} est un segment qui relie un \solution[3cm]{sommet} au \solution[2.5cm]{milieu} du côté opposé.

	Un triangle possède \solution[1.5cm]{trois} médianes.
\end{definitionbox}

\begin{examplebox}
	\textbf{Les trois médianes du triangle $ABC$}

	\ifcorrige
	\begin{center}
		\begin{tikzpicture}[scale=0.9]
			% Triangle
			\coordinate (A) at (0,0);
			\coordinate (B) at (6,0);
			\coordinate (C) at (2,4);

			% Milieux
			\coordinate (I) at ($(B)!0.5!(C)$);
			\coordinate (J) at ($(A)!0.5!(C)$);
			\coordinate (K) at (3,0);

			% Triangle
			\draw[thick] (A) -- (B) -- (C) -- cycle;

			% Médianes
			\draw[mainBlue, thick] (A) -- (I);
			\draw[mainBlue, thick] (B) -- (J);
			\draw[mainBlue, thick] (C) -- (K);

			% Centre de gravité
			\coordinate (G) at (barycentric cs:A=1,B=1,C=1);
			\fill[mainGreen] (G) circle (3pt) node[below right] {$G$};

			% Points
			\fill (A) circle (2pt) node[below left] {$A$};
			\fill (B) circle (2pt) node[below right] {$B$};
			\fill (C) circle (2pt) node[above] {$C$};
			\fill (I) circle (2pt) node[right] {$I$};
			\fill (J) circle (2pt) node[left] {$J$};
			\fill (K) circle (2pt) node[below] {$K$};

			% Légendes
			\node[mainBlue] at (4.5,2.5) {médiane $[AI]$};
			\node[mainBlue] at (-0.5,2.5) {médiane $[BJ]$};
			\node[mainBlue] at (2.5,1.5) {médiane $[CK]$};
		\end{tikzpicture}
	\end{center}
	\fi

	\begin{itemize}
		\item $I$ est le milieu de \solution[2cm]{$[BC]$}, donc $[AI]$ est une médiane
		\item $J$ est le milieu de \solution[2cm]{$[AC]$}, donc $[BJ]$ est une médiane
		\item $K$ est le milieu de \solution[2cm]{$[AB]$}, donc $[CK]$ est une médiane
	\end{itemize}
\end{examplebox}

\begin{proprietebox}[Centre de gravité]
	Les trois médianes d'un triangle sont \solution[4cm]{concourantes}.

	Le point de concours des trois médianes s'appelle le \textbf{\solution[5cm]{centre de gravité}} (ou \textbf{barycentre}) du triangle. On le note souvent $G$.
\end{proprietebox}

\begin{proprietebox}[Propriété du centre de gravité]
	Le centre de gravité $G$ se situe aux \solution[3cm]{deux tiers} de chaque médiane en partant du sommet.

	Par exemple, si $I$ est le milieu de $[BC]$, alors :
	$$AG = \solution[2cm]{\frac{2}{3}} \times AI \quad \text{et} \quad GI = \solution[2cm]{\frac{1}{3}} \times AI$$
\end{proprietebox}

\begin{remarkbox}
	\textbf{Propriété physique :}

	Si on découpe un triangle dans du carton, le centre de gravité est le point d'équilibre du triangle. En plaçant un doigt sous le centre de gravité, le triangle reste en équilibre.
\end{remarkbox}

\begin{methodebox}[Construire une médiane]
	Pour construire la médiane issue de $A$ dans le triangle $ABC$ :
	\begin{enumerate}
		\item Avec le \solution[2.5cm]{compas}, trouver le milieu $I$ du segment $[BC]$.
		\item Tracer le segment \solution[2cm]{$[AI]$}.
		\item Ce segment est la médiane issue de \solution[1.5cm]{$A$}.
	\end{enumerate}
\end{methodebox}

\subsection{Les médiatrices}

\begin{definitionbox}[Médiatrice d'un segment]
	La \textbf{médiatrice} d'un segment est la droite qui passe par le \solution[2.5cm]{milieu} du segment et qui est \solution[4cm]{perpendiculaire} à ce segment.
\end{definitionbox}

\begin{proprietebox}[Propriété caractéristique de la médiatrice]
	Un point $M$ appartient à la médiatrice du segment $[AB]$ si et seulement si :
	$$MA = \solution[1.5cm]{MB}$$

	\textbf{Autrement dit :} La médiatrice d'un segment est l'ensemble des points \solution[5cm]{équidistants} des extrémités du segment.
\end{proprietebox}

\begin{definitionbox}[Médiatrices d'un triangle]
	Dans un triangle $ABC$, on peut tracer les \textbf{médiatrices} des trois côtés :
	\begin{itemize}
		\item La médiatrice du côté $[AB]$
		\item La médiatrice du côté $[BC]$
		\item La médiatrice du côté $[AC]$
	\end{itemize}
\end{definitionbox}

\begin{examplebox}
	\textbf{Les trois médiatrices du triangle $ABC$}

	\ifcorrige
	\begin{center}
		\begin{tikzpicture}[scale=0.9]
			% Triangle
			\coordinate (A) at (0,0);
			\coordinate (B) at (6,0);
			\coordinate (C) at (2,4);

			% Milieux
			\coordinate (I) at ($(B)!0.5!(C)$);
			\coordinate (J) at ($(A)!0.5!(C)$);
			\coordinate (K) at (3,0);

			% Centre du cercle circonscrit
			\coordinate (O) at (3,1.5);

			% Triangle
			\draw[thick] (A) -- (B) -- (C) -- cycle;

			% Médiatrices
			\draw[mainRed, thick] (K) -- ++(0,4.5);
			\draw[mainRed, thick] (I) -- ++(-1.5,1) -- ++(3,-2);
			\draw[mainRed, thick] (J) -- ++(1.5,1) -- ++(-3,2);

			% Cercle circonscrit
			\draw[mainGreen, dashed] (O) circle (2.5);

			% Centre du cercle circonscrit
			\fill[mainBlue] (O) circle (3pt) node[above right] {$O$};

			% Marques de milieux
			\fill (I) circle (2pt);
			\fill (J) circle (2pt);
			\fill (K) circle (2pt);

			% Marques d'angles droits
			\draw (K) ++(0,0.3) -- ++(0.3,0) -- ++(0,-0.3);

			% Points
			\fill (A) circle (2pt) node[below left] {$A$};
			\fill (B) circle (2pt) node[below right] {$B$};
			\fill (C) circle (2pt) node[above] {$C$};
		\end{tikzpicture}
	\end{center}
	\fi
\end{examplebox}

\begin{proprietebox}[Centre du cercle circonscrit]
	Les trois médiatrices d'un triangle sont \solution[4cm]{concourantes}.

	Le point de concours des trois médiatrices s'appelle le \textbf{\solution[7cm]{centre du cercle circonscrit}} au triangle. On le note souvent $O$.
\end{proprietebox}

\begin{proprietebox}[Cercle circonscrit]
	Le centre $O$ du cercle circonscrit est \solution[5cm]{équidistant} des trois sommets du triangle :
	$$OA = OB = OC$$

	Le cercle de centre $O$ et de rayon $OA$ passe par les \solution[2cm]{trois} sommets du triangle. On l'appelle le \textbf{cercle circonscrit} au triangle.
\end{proprietebox}

\begin{remarkbox}
	\textbf{Cas particuliers :}
	\begin{itemize}
		\item Dans un triangle rectangle, le centre du cercle circonscrit est le milieu de l'hypoténuse.
		\item Dans un triangle obtusangle, le centre du cercle circonscrit est à l'extérieur du triangle.
	\end{itemize}
\end{remarkbox}

\begin{methodebox}[Construire une médiatrice]
	Pour construire la médiatrice du segment $[AB]$ :

	\textbf{Méthode 1 : Avec le compas}
	\begin{enumerate}
		\item Tracer un arc de cercle de centre $A$ et de rayon supérieur à la moitié de $AB$.
		\item Tracer un arc de cercle de centre $B$ et de même rayon.
		\item Les deux arcs se coupent en deux points. La droite passant par ces deux points est la médiatrice.
	\end{enumerate}

	\textbf{Méthode 2 : Avec la règle et l'équerre}
	\begin{enumerate}
		\item Trouver le milieu $I$ du segment $[AB]$ avec le compas.
		\item Tracer la perpendiculaire à $(AB)$ passant par $I$ avec l'équerre.
	\end{enumerate}
\end{methodebox}

\subsection{Tableau récapitulatif}

\begin{center}
	\begin{tabular}{|p{3cm}|p{4.5cm}|p{4.5cm}|p{3cm}|}
		\hline
		\textbf{Droite remarquable} & \textbf{Définition} & \textbf{Construction} & \textbf{Point remarquable} \\
		\hline
		\textbf{Hauteur} & Droite passant par un sommet et perpendiculaire au côté opposé & Perpendiculaire avec équerre ou compas & Orthocentre (point de concours des hauteurs) \\
		\hline
		\textbf{Médiane} & Segment joignant un sommet au milieu du côté opposé & Trouver le milieu puis tracer & Centre de gravité (point de concours des médianes) \\
		\hline
		\textbf{Médiatrice} & Droite perpendiculaire à un côté passant par son milieu & Arcs de cercle de même rayon & Centre du cercle circonscrit (point de concours des médiatrices) \\
		\hline
	\end{tabular}
\end{center}

\begin{astucebox}
	\textbf{Pour ne pas confondre :}
	\begin{itemize}
		\item \textbf{Hauteur} : passe par un \textbf{sommet}, perpendiculaire au côté opposé
		\item \textbf{Médiane} : passe par un \textbf{sommet} et le \textbf{milieu} du côté opposé
		\item \textbf{Médiatrice} : passe par le \textbf{milieu} d'un côté, perpendiculaire à ce côté (ne passe pas nécessairement par un sommet)
	\end{itemize}
\end{astucebox}

\begin{quizbox}
	\textbf{Auto-évaluation}
	\begin{enumerate}[label=\arabic*)]
		\item Énoncer l'inégalité triangulaire.
		\item Trois longueurs 4 cm, 7 cm et 12 cm peuvent-elles former un triangle ?
		\item Qu'est-ce qu'une hauteur d'un triangle ?
		\item Comment s'appelle le point de concours des médianes ?
		\item Quelle propriété caractérise la médiatrice d'un segment ?
	\end{enumerate}
\end{quizbox}

\begin{summarybox}
	\textbf{L'essentiel à retenir}
	\begin{itemize}
		\item \textbf{Inégalité triangulaire :} Dans un triangle, la longueur d'un côté est toujours strictement inférieure à la somme des deux autres.
		\item Pour qu'un triangle existe, il faut que la plus grande longueur soit strictement inférieure à la somme des deux autres.
		\item Un triangle possède 3 hauteurs qui se coupent en l'orthocentre.
		\item Un triangle possède 3 médianes qui se coupent au centre de gravité (à 2/3 de chaque médiane depuis le sommet).
		\item Un triangle possède 3 médiatrices qui se coupent au centre du cercle circonscrit.
		\item Le cercle circonscrit passe par les trois sommets du triangle.
	\end{itemize}
\end{summarybox}

\section{Exercices d'application}

\begin{exercisebox}
	\textbf{Exercice 1 : Inégalité triangulaire} ($\star$ Exercice de base)

	Pour chacune des propositions suivantes, dire si on peut construire un triangle avec les longueurs données. Justifier.

	\begin{multicols}{2}
		\begin{enumerate}[label=\alph*)]
			\item 5 cm, 7 cm et 9 cm
			\item 3 cm, 8 cm et 12 cm
			\item 6 cm, 6 cm et 6 cm
			\item 2 cm, 5 cm et 7 cm
			\item 10 cm, 4 cm et 5 cm
			\item 8 cm, 15 cm et 20 cm
		\end{enumerate}
	\end{multicols}

	\medskip

	\textbf{Exercice 2 : Trouver des longueurs possibles}

	Un triangle $ABC$ a pour côtés $AB = 8$ cm et $BC = 5$ cm.

	\begin{enumerate}[label=\alph*)]
		\item Donner un encadrement de la longueur $AC$.
		\item Le côté $[AC]$ peut-il mesurer 3 cm ? 4 cm ? 13 cm ? 10 cm ?
		\item Donner trois valeurs possibles pour la longueur $AC$.
	\end{enumerate}

	\medskip

	\textbf{Exercice 3 : Constructions}

	\begin{enumerate}[label=\alph*)]
		\item Construire le triangle $DEF$ tel que $DE = 7$ cm, $EF = 5$ cm et $DF = 6$ cm.
		\item Construire le triangle $GHI$ tel que $GH = 6$ cm, $GI = 5$ cm et $\widehat{IGH} = 50^\circ$.
		\item Construire le triangle $JKL$ tel que $JK = 7$ cm, $\widehat{KJL} = 40^\circ$ et $\widehat{JKL} = 60^\circ$.
	\end{enumerate}
\end{exercisebox}

\begin{exercisebox}
	\textbf{Exercice 4 : Hauteurs} ($\star$ Exercice de base)

	\begin{enumerate}[label=\alph*)]
		\item Construire un triangle $ABC$ quelconque.
		\item Tracer la hauteur issue de $A$.
		\item Tracer la hauteur issue de $B$.
		\item Tracer la hauteur issue de $C$.
		\item Marquer le point $H$, orthocentre du triangle.
		\item Vérifier que les trois hauteurs passent bien par le point $H$.
	\end{enumerate}

	\medskip

	\textbf{Exercice 5 : Médianes}

	\begin{enumerate}[label=\alph*)]
		\item Construire un triangle $DEF$ avec $DE = 8$ cm, $EF = 6$ cm et $DF = 7$ cm.
		\item Placer les milieux $I$, $J$ et $K$ des côtés $[EF]$, $[DF]$ et $[DE]$.
		\item Tracer les trois médianes $[DI]$, $[EJ]$ et $[FK]$.
		\item Marquer le point $G$, centre de gravité du triangle.
		\item Mesurer les longueurs $DG$ et $GI$. Vérifier que $DG = 2 \times GI$.
	\end{enumerate}

	\medskip

	\textbf{Exercice 6 : Médiatrices}

	\begin{enumerate}[label=\alph*)]
		\item Construire un triangle $RST$ quelconque.
		\item Construire la médiatrice du côté $[RS]$.
		\item Construire la médiatrice du côté $[ST]$.
		\item Construire la médiatrice du côté $[RT]$.
		\item Marquer le point $O$, centre du cercle circonscrit.
		\item Tracer le cercle de centre $O$ passant par le point $R$. Vérifier qu'il passe aussi par $S$ et $T$.
	\end{enumerate}
\end{exercisebox}

\begin{exercisebox}
	\textbf{Exercice 7 : Triangle rectangle} ($\star\star$ Approfondissement)

	Construire un triangle $ABC$ rectangle en $A$ tel que $AB = 6$ cm et $AC = 4$ cm.

	\begin{enumerate}[label=\alph*)]
		\item Tracer les trois hauteurs. Que remarque-t-on ?
		\item Où se situe l'orthocentre ?
		\item Tracer les trois médiatrices. Où se situe le centre du cercle circonscrit ?
		\item Quelle est la particularité du cercle circonscrit à un triangle rectangle ?
	\end{enumerate}

	\medskip

	\textbf{Exercice 8 : Problèmes}

	\begin{enumerate}[label=\alph*)]
		\item Le centre de gravité $G$ d'un triangle $ABC$ se situe sur la médiane issue de $A$ à 6 cm du sommet $A$. Quelle est la longueur totale de cette médiane ?

		\item Un triangle $ABC$ a un côté $[AB]$ de longueur 12 cm. Le point $M$ est le milieu de $[AB]$. Le centre du cercle circonscrit $O$ se trouve sur la médiatrice de $[AB]$. Si $OA = 10$ cm, calculer $OM$.

		\item Dans un triangle isocèle $ABC$ avec $AB = AC$, montrer que la médiane, la hauteur et la médiatrice issues de $A$ sont confondues.
	\end{enumerate}
\end{exercisebox}

\begin{exercisebox}[Remédiation] ($\star$ Pour les élèves en difficulté)
	\textbf{Exercice R1 :} Vocabulaire

	Compléter les phrases suivantes :
	\begin{enumerate}[label=\alph*)]
		\item Une hauteur est une droite passant par un \solution[3cm]{sommet} et \solution[5cm]{perpendiculaire} au côté opposé.
		\item Une médiane relie un sommet au \solution[2.5cm]{milieu} du côté opposé.
		\item Une médiatrice est perpendiculaire à un côté et passe par son \solution[2.5cm]{milieu}.
	\end{enumerate}

	\medskip

	\textbf{Exercice R2 :} Inégalité triangulaire simple

	On donne deux côtés d'un triangle : 5 cm et 7 cm. Le troisième côté peut-il mesurer :
	\begin{enumerate}[label=\alph*)]
		\item 1 cm ? \solution[2cm]{Non}
		\item 3 cm ? \solution[2cm]{Oui}
		\item 12 cm ? \solution[2cm]{Oui}
		\item 13 cm ? \solution[2cm]{Non}
	\end{enumerate}

	\medskip

	\textbf{Exercice R3 :} Construction guidée

	\begin{enumerate}[label=\alph*)]
		\item Construire un triangle $ABC$ avec $AB = 6$ cm, $AC = 5$ cm et $BC = 4$ cm.
		\item Placer le milieu $I$ du segment $[BC]$.
		\item Tracer la médiane $[AI]$.
	\end{enumerate}
\end{exercisebox}

\begin{exercisebox}[Approfondissement] ($\star\star\star$ Pour aller plus loin)
	\textbf{Exercice A1 :} Droite d'Euler

	Dans un triangle quelconque (non équilatéral), construire :
	\begin{enumerate}[label=\alph*)]
		\item L'orthocentre $H$
		\item Le centre de gravité $G$
		\item Le centre du cercle circonscrit $O$
	\end{enumerate}

	Que remarque-t-on concernant l'alignement de ces trois points ?

	\solution[8cm]{Ces trois points sont alignés sur une droite appelée « droite d'Euler ». De plus, $G$ est situé aux 2/3 du segment $[OH]$ en partant de $O$.}

	\medskip

	\textbf{Exercice A2 :} Cercle d'Euler (cercle des neuf points)

	Dans un triangle $ABC$, on considère les neuf points suivants :
	\begin{itemize}
		\item Les trois milieux des côtés
		\item Les trois pieds des hauteurs
		\item Les trois milieux des segments joignant les sommets à l'orthocentre
	\end{itemize}

	Construire ces neuf points. Que remarque-t-on ?

	\solution[8cm]{Ces neuf points sont situés sur un même cercle, appelé « cercle d'Euler » ou « cercle des neuf points ».}

	\medskip

	\textbf{Exercice A3 :} Triangle équilatéral

	Dans un triangle équilatéral, montrer que :
	\begin{enumerate}[label=\alph*)]
		\item Les hauteurs, les médianes et les médiatrices sont confondues
		\item L'orthocentre, le centre de gravité et le centre du cercle circonscrit sont confondus
	\end{enumerate}

	\medskip

	\textbf{Exercice A4 :} Optimisation

	On veut construire un pont reliant trois villages $A$, $B$ et $C$ formant un triangle. On souhaite placer un point $P$ tel que la somme des distances $PA + PB + PC$ soit minimale.

	Quel point remarquable du triangle doit-on choisir pour $P$ ?

	\solution[8cm]{On doit choisir le point de Fermat (point de Torricelli), qui minimise la somme des distances aux trois sommets. Dans un triangle dont tous les angles sont inférieurs à 120°, ce point est tel que les angles $\widehat{APB}$, $\widehat{BPC}$ et $\widehat{CPA}$ valent chacun 120°.}
\end{exercisebox}
